\begin{center}\textls[100]{c. Kubische und biquadratische Reste. Komplexe ganze Zahlen.}\end{center}
\begin{center}\bfseries 21.\end{center}
\noindent\hspace*{\parindent}Wie gesagt, sind die letzten Beweise des Reziprozitätsgesetzes zumeist durch die Untersuchungen hervorgerufen, welche \textsc{Gauss} über kubische und biquadratische Reste unternahm, Untersuchungen, zu denen überzugehen für ihn nahe lag, nachdem er die Theorie der quadratischen Reste erledigt hatte. Er sagt darüber in einem Briefe an \textsc{Dirichlet} vom 30. Mai 1828 (W. II, S. 516), was folgt: »Die ganze Untersuchung, deren Stoff ich schon seit 23 Jahren vollständig besitze, die Beweise der Haupttheoreme aber (zu welchen das in der ersten Commentation noch nicht zu rechnen ist) seit etwa 14 Jahren« usw. — Fällt hiernach der Beginn seiner bezüglichen Forschungen spätestens in das Jahr 1805, so wird dieses Jahr von \textsc{Gauss} auch früher schon, wo seine Erinnerung sicher zuverlässiger war, zu wiederholten Malen (s. W. II, S. 50 und 161, \textit{Comm. prima}, W. II, S. 67 und S. 165 in der Anzeige derselben, \textit{Comm. secunda} W. II, S. 102) bestimmt als deren Anfang angegeben. Dem braucht die Tagebuchnotiz Nr. 130 vom 15. Februar 1807: »theoria residuorum cubicorum et biquadraticorum incepta« nicht zu widersprechen, da man in ihr den Anfang einer planmässig geordneten Entwicklung oder Darstellung seiner Ergebnisse erblicken darf, welche nach den Nummern 132, 133 T. und nach \textsc{Gauss}’ am 30. April 1807 an \textsc{Sophie Germain} gerichteten Briefe (abgedruckt W. X1, S. 70) schon eine gewisse Höhe erreicht hatten. Diese Ergebnisse finden sich dann schliesslich in zwei Abhandlungen, der \textit{Theoria residuorum biquadraticorum, Comm. I u. II.} (Comm. Gotting. rec. \textbf{6}, 1828, vorgelegt am 5. April 1828 und \textbf{7}, 1832, vorgelegt am 15. April 1831), welche \textsc{Gauss} als näher verwandt mit derjenigen der quadratischen Reste einer Darstellung der kubischen vor-
\newpage\SourceLeaf{60}\ArithmeticHead{53}{R}{KUBISCHE UND BIQUADRATISCHE RESTE. KOMPLEXE GANZE ZAHLEN.}
\noindent zog, wenn auch leider nur zu einem Teile veröffentlicht. Sie müssen zunächst in raschem Fortschritte gelungen sein, wie die Notizen Nr. 131—133 T. vom 17., 22., 24. Februar bezeugen, die ohne Zweifel die Auffindung der wesentlichsten Sätze der \textit{Comm. prima} bedeuten. Dies sind in der Hauptsache die folgenden.

Beschränkt man sich auf den Fall der Primzahlen \(p=4n+1\), welcher einzig Schwierigkeiten verursacht, so zerfallen die modulo \(p\) inkongruenten Zahlen \(z\), welche sich bisher nur in zwei Klassen, quadratische Reste und Nichtreste, unterschieden, jenachdem \(z^{\frac{p-1}{2}}\equiv+1\) oder \(-1\) (mod. \(p\)) war, in der Theorie der biquadratischen Reste in vier Klassen \(A,B,C,D\) von je \(\frac{p-1}{4}\) Zahlen, die bezüglich den Kongruenzen
\[
z^{\frac{p-1}{4}}\equiv1,f,f^2,f^3\ (\mathrm{mod.}\ p)
\]
genügen, in denen \(f^2\equiv-1\) (mod. \(p\)); die Klassen \(A,C\) zusammen sind die quadratischen Reste, \(B,D\) zusammen die quadratischen Nichtreste.

Die \textit{Comm. prima} lehrt nun, in welche dieser Klassen die Zahlen \(-1\) und \(2\) gehören. Die erstere gehört zu \(A\) oder \(C\), jenachdem \(p\equiv1\) oder \(5\) (mod. 8). Die Zahl \(2\) ist quadratischer Nichtrest bezüglich der Moduln \(p\equiv5\) (mod. 8). Ist aber \(p\equiv1\) (mod. 8), so liess die Induktion \textsc{Gauss} einen Zusammenhang der Frage mit der Zerlegung \(p=\alpha^2+2\beta^2\) erkennen, wonach \(2\) zu \(A\) oder \(C\) gehört, jenachdem \(\alpha\equiv\pm1\) oder \(\pm3\) (mod 8) ist, ein Satz, der unschwer zu begründen war.

Hierdurch fand sich nun \textsc{Gauss} zu der Untersuchung geführt, ob ein ähnlicher Zusammenhang auch mit der für jede Primzahl \(p=4n+1\) möglichen Zerlegung \(p=a^2+b^2\) stattfinde. Das ist in der Tat der Fall, aber er liegt tiefer. Die gedachte Zerlegung ist zunächst eng mit den Zahlen \(m_k^{(h)}\) verbunden, welche bei der Bildung der biquadratischen Gleichung für die vier \(\frac{p-1}{4}\)-gliedrigen Perioden \(p^{\mathrm{ter}}\) Einheitswurzeln auftreten und die Anzahl der Lösungen \(t,u\) der Kongruenz
\[
1+g^{4t+h}\equiv g^{4u+k}\ (\mathrm{mod.}\ p)
\]
bezeichnen, wo \(g\) eine primitive Wurzel (mod. \(p\)) ist. Die Beziehungen zwischen ihnen, sowie die Werte von \(a,b\) bei der durch sie gelieferten Zerlegung
\newpage\SourceLeaf{61}\ArithmeticHead{54}{L}{BACHMANN, ÜBER GAUSS’ ZAHLENTHEORETISCHE ARBEITEN.}
\noindent \(p=a^2+b^2\) ergeben den Satz, dass \(2\) beziehungsweise zu \(A,B,C,D\) gehört, jenachdem
\[
\frac{b}{2}\equiv0,1,2,3\ (\mathrm{mod.}\ 4).
\]
Übrigens hängt das Vorzeichen von \(b\) und somit die Zugehörigkeit von \(2\) zu einer der Klassen \(B,D\) von der Wahl der Zahl \(f\) unter den beiden Wurzeln der Kongruenz \(f^2\equiv-1\), d. i. von der Wahl der primitiven Wurzel \(g\) ab, und es besteht zwischen \(a,b\) und \(f\) die Kongruenz \(b\equiv af\) (mod. \(p\)), eine Tatsache, die \textsc{Gauss} (Nr. 133 T.) am 24. Februar 1807 gefunden zu haben scheint, und welche verstattet, den biquadratischen Charakter der Zahl \(2\) auch durch die Kongruenz
\[
b'^{\frac{a'b'}{2}}\equiv a'^{\frac{a'b'}{2}}\mathbin{.}2^{\frac{p-1}{4}}\ (\mathrm{mod.}\ p),
\]
wo \(a',b'\) in der Formel \(p=a'^2+b'^2\) positiv, \(a'\) ungerade gedacht sind, auszusprechen (W. II, S. 96). Nebenbei finden sich \(a,b\) mittels der Entwicklung von
\[
(x^4+1)^{\frac{p-1}{4}}\ (\mathrm{mod.}\ p)
\]
als die absolut kleinsten Reste modulo \(p\), welche den Kongruenzen
\[
2a\equiv\frac{r}{s},\quad 2b\equiv\pm r^2\ (\mathrm{mod.}\ p)
\]
mit
\[
s=1\mathbin{.}2\mathbin{.}3\ldots\frac{p-1}{4},\quad r=\frac{p+3}{4}\mathbin{.}\frac{p+7}{4}\ldots\frac{p-1}{2}
\]
genügen. Ihnen zufolge ist \(a\equiv1\) (mod. 4) und hierdurch seinem Vorzeichen nach bestimmt, während über dasjenige von \(b\) sich ohne weiteres nicht entscheiden lässt (s. unten S. 69). Ähnliche Sätze hat \textsc{Gauss} noch mehrere gegeben in einer Notiz des Nachlasses, die W. X1, S. 39 abgedruckt ist.
